% Generated draft: requires mathematical and source review.
\documentclass[11pt]{article}
\usepackage[margin=25mm]{geometry}
\usepackage{fontspec}
\setmainfont{Latin Modern Roman}
\usepackage{amsmath,amssymb,mathtools,unicode-math}
\setmathfont{Latin Modern Math}
\usepackage{enumitem,needspace,xurl,hyperref}
\setlength{\parindent}{0pt}
\setlength{\parskip}{5pt}
\setlength{\emergencystretch}{3em}
\newcommand{\N}{\mathbb{N}}
\newcommand{\Z}{\mathbb{Z}}
\newcommand{\Q}{\mathbb{Q}}
\newcommand{\R}{\mathbb{R}}
\newcommand{\C}{\mathbb{C}}
\newcommand{\F}{\mathbb{F}}
\newcommand{\A}{\mathbb{A}}
\newcommand{\PP}{\mathbb{P}}
\newcommand{\E}{\mathbb{E}}
\newcommand{\eps}{\varepsilon}
\newcommand{\dd}{\,\mathrm{d}}
\begin{document}
\section*{UnsolvedMath ID: 11000234}
\textbf{Source problem code:} catalogue entry
\subsection{Definitions and mathematical statement}
Let $S$ be the surface in the source and let $\operatorname{MCG}(S)$ act on its Teichm\"uller space with the Teichm\"uller metric. In the customary mapping-class-group usage, a Schottky subgroup is a finitely generated free subgroup that is convex cocompact: a Teichm\"uller orbit is quasiconvex, meaning that geodesic segments joining its points stay in a uniform neighborhood of that orbit. The source's exact rank and low-complexity conventions should be confirmed.

A subgroup $H\le\operatorname{MCG}(S)$ is \emph{virtual Schottky} if it has a finite-index subgroup $F\le H$ which is Schottky. This definition is explicitly supplied in the adjacent source item, local ID 11000233. Thus the required property involves both a finite-index free subgroup and its geometric convex-cocompactness, not just abstract virtual freeness or the pseudo-Anosov nature of individual elements.

\textbf{Problem 3.6.} Give examples and constructions of virtual Schottky subgroups of $\operatorname{MCG}(S)$. For a proposed construction, identify the subgroup $H$, exhibit or characterize its finite-index Schottky subgroup $F$, and establish the required geometric property. The question does not ask that every abstract virtually free group embed this way, nor explicitly require that $H$ itself be nonfree.

The source announces a construction due to Honglin Min, described at the time as a doctoral candidate at Rutgers University, Newark. The excerpt ends at the colon introducing that construction, so its actual data and theorem are absent. No construction is reconstructed or attributed to Min without those missing details.
\subsection{Short English statement}
Construct mapping-class subgroups with finite-index free convex-cocompact subgroups. The source announces Honglin Min's example, but the construction itself is not included.
\subsection{Sources}
\begin{itemize}
\item Local UnsolvedMath registry, permanent ID 11000234, source code catalogue entry; no separate proposer or proposal year supplied.
\item Farb, edited volume \emph{Problems on Mapping Class Groups and Related Topics}, Problem 3.6, PDF page 272; \url{https://www.math.uchicago.edu/~farb/papers/mcgbook.pdf}. The registry records PDF-font-structure extraction and access on 2026-07-29; the inventory distinguishes 329 labeled headings by source page.
\item Adjacent local ID 11000233 gives the finite-index Schottky definition. The present source credits Honglin Min with a following construction, without reproducing it or giving a separate bibliographic citation.
\end{itemize}
\end{document}
